\section{Conclusion}
\label{sec:conclusion}

We opened by observing that HumanEval-style algorithm training outperforms MBPP utility-script
training on MBPP's own benchmark --- across all three random seeds. After tracing the causal
mechanism through embedding-space distributional alignment, we can now explain why: the model
trained on the data most similar (in code-embedding space) to the test distribution achieves
the highest performance, regardless of whether surface-level benchmark names suggest otherwise.

In this work, we:
\begin{enumerate}
\item Showed that SFT source identity causes up to 31.9pp performance differences under controlled conditions (F=11.37, p=0.020).
\item Demonstrated that embedding alignment perfectly predicts the performance rank across four conditions ($\rho$=1.0, p=0.042).
\item Discovered that algorithmic function-completion training generalizes better than utility-script training across benchmarks.
\item Identified $\eta^2$ as an inappropriate metric for cross-scale comparison when seed variance changes with scale.
\end{enumerate}

\textbf{Future work} includes: completing the full MBPP+ 4-condition evaluation on existing checkpoints;
using embedding alignment as an active SFT data selection criterion; testing whether the alignment
mechanism holds across additional coding languages and benchmarks; and understanding whether the
HumanEval advantage holds with the MBPP native prompt format.

Understanding what code SFT actually teaches --- which distributional properties of training data
transfer to test benchmarks, and why --- is essential for reproducible benchmark evaluation and
principled data curation. We hope this work, by providing both controlled empirical evidence and
the embedding alignment framework, accelerates that understanding.
